\documentclass[11pt,a4paper]{article}
\usepackage[T1]{fontenc}
\usepackage{lmodern}
\usepackage[margin=25mm]{geometry}
\usepackage{amsmath,amssymb,amsthm,mathtools}
\usepackage{enumitem,microtype,xcolor,booktabs,tabularx}
\usepackage[hyperpageref]{backref}
\usepackage[pagebackref]{hyperref}
\hypersetup{
    colorlinks=true,
    linkcolor=red,
    filecolor=blue,
    urlcolor=orange,
    citecolor=violet
}
\renewcommand*{\backref}[1]{}
\renewcommand*{\backrefalt}[4]{
  \ifcase #1 Not cited.
  \or Cited on page #2.
  \else Cited on pages #2.
  \fi}
\definecolor{darkblue}{rgb}{0,0,0.7}
\newcommand{\defn}[1]{\textsl{\color{darkblue} #1}}
\setlist[enumerate]{leftmargin=*,itemsep=.5em,topsep=.4em}
\setlength{\emergencystretch}{2em}
\newcommand{\D}{\mathcal D}
\renewcommand{\H}{\mathcal H}
\newcommand{\ZZ}{\mathbb{Z}}
\newcommand{\op}{\mathrm{op}}
\newcommand{\RR}{\mathbf{R}}
\newcommand{\std}{\Delta}
\newcommand{\costd}{\nabla}
\newcommand{\dual}{\mathbb D}
\newcommand{\Ker}{\operatorname{Ker}}
\DeclareMathOperator{\fd}{fd}
\DeclareMathOperator{\Hom}{Hom}
\DeclareMathOperator{\End}{End}
\DeclareMathOperator{\REnd}{\RR End}
\DeclareMathOperator{\add}{add}
\DeclareMathOperator{\thick}{thick}
\DeclareMathOperator{\per}{per}
\DeclareMathOperator{\Filt}{Filt}
\let\mod\relax
\DeclareMathOperator{\mod}{mod}
\theoremstyle{plain}
\newtheorem{theorem}{Theorem}[section]
\newtheorem{proposition}[theorem]{Proposition}
\newtheorem{corollary}[theorem]{Corollary}
\theoremstyle{definition}
\newtheorem{definition}[theorem]{Definition}
\newtheorem{example}[theorem]{Example}
\numberwithin{equation}{section}
\title{$m$-qh dgas}
\author{}
\date{}

\begin{document}
\maketitle
\begin{abstract}
We formulate quasi-heredity for basic proper connective dg algebras
using exceptional standard objects and projective presentations in
$\H_A^m$. Kernels admit finite filtrations by non-negative shifts of
standards with larger weights, with every factor in $\H_A^m$.
The axioms imply fullness in $\D_{\fd}(A)=\per A$, determine the
standard objects from the ordered projectives, and place the left-dual
objects in $\H_A^m[1-m]$. Graded path dg algebras admit such structures
for every ordering at a width controlled by path degrees. We obtain
an intrinsic criterion for shifted filtrations, an exact least-width
formula, and the corresponding injective presentations. The remaining
foundational questions concern recognition and comparison of the axioms.
We compare replacements for the exceptional-sequence axiom, with
counterexamples and a separate summary for each proposal. Finite
projective tests can make exceptionality and smoothness consequences;
tests using only non-negative shifts fail to control negative cohomology.
\end{abstract}

\section{Conventions and definition}\label{sec:definition}

Fix an algebraically closed field $\Bbbk$ and integers $m,n\geq1$.
All dg modules are right modules, and complexes are cochain complexes:
$d_X:X^q\to X^{q+1}$ and $H^q(X[s])=H^{q+s}(X)$.
Let $A$ be a proper connective dg algebra with
\begin{equation}\label{eq:standing}
 \dim_{\Bbbk}H^*(A)<\infty,\qquad
 H^q(A)=0\quad(q\notin[1-m,0]).
\end{equation}
Assume that $B=H^0(A)$ is basic and that closed degree-zero orthogonal
idempotents $e_1,\ldots,e_n$ sum to $1$, with classes forming a complete
set of primitive idempotents of $B$. Use $e_v$ also for its class in $B$,
and let $S_v$ be the simple top of $e_vB$. Set
\begin{equation}\label{eq:categories}
 \begin{gathered}
 P_v=e_vA,\qquad
 \D_{\fd}(A)=\{X\in\D(A):\dim_{\Bbbk}H^*(X)<\infty\},\\
 \per A=\thick_{\D(A)}(A),\qquad
 \H_A^m=\D_{\fd}(A)^{[1-m,0]}.
 \end{gathered}
\end{equation}
The notation $\thick$ denotes closure under shifts, cones, and direct
summands. Properness gives $\per A\subseteq\D_{\fd}(A)$.
For a complex $V$ of vector spaces, write
$\dual V=\Hom^\bullet_{\Bbbk}(V,\Bbbk)$ for its complex dual.

The category $\H_A^m$ is extriangulated, with
\begin{equation}\label{eq:extriangulated}
 \begin{gathered}
 \mathbb E_{\H_A^m}(X,Y)=\Hom_{\D_{\fd}(A)}(X,Y[1]),\\
 \mathcal P(\H_A^m)=\add A,\qquad
 \mathcal I(\H_A^m)=\add(\dual A[m-1]).
 \end{gathered}
\end{equation}
Here $X,Y\in\H_A^m$; see \cite[Corollary 2.3]{MP25}.
A conflation $K\to P\to X\dashrightarrow$ is induced by a triangle
$K\to P\to X\to K[1]$ with all three objects in $\H_A^m$.
For $\mathcal S\subseteq\H_A^m$, write $\Filt_{\H_A^m}(\mathcal S)$
for the objects $M$ admitting a finite sequence of conflations
\begin{equation}\label{eq:filtration}
 M_{a-1}\longrightarrow M_a\longrightarrow F_a\dashrightarrow,
 \quad M_0=0,\quad M_\ell=M,\quad F_a\in\mathcal S
 \quad(1\leq a\leq\ell).
\end{equation}
All $M_a$ belong to $\H_A^m$. Empty filtrations are allowed, and finite
direct sums belong to $\Filt_{\H_A^m}(\mathcal S)$; closure under direct
summands is not imposed. The inflations in \eqref{eq:filtration} need
not be monomorphisms of dg modules.

An object $E\in\D_{\fd}(A)$ is \defn{exceptional} if
\begin{equation}\label{eq:exceptional}
 \Hom_{\D_{\fd}(A)}(E,E[t])\cong
 \begin{cases}\Bbbk,&t=0,\\0,&t\neq0\end{cases}
\qquad(t\in\ZZ).
\end{equation}
An ordered tuple $(E_1,\ldots,E_n)$ of exceptional objects is an
\defn{exceptional sequence} if
\begin{equation}\label{eq:exceptional-sequence}
 \Hom_{\D_{\fd}(A)}(E_i,E_j[t])=0
 \qquad(1\leq j<i\leq n,\ t\in\ZZ).
\end{equation}
The indices $i,j$ in \eqref{eq:exceptional-sequence} denote positions
in the sequence.
For exceptional objects $\std_1,\ldots,\std_n$, a permutation
$\sigma\in S_n$ lists the vertex labels in the standard sequence
\begin{equation}\label{eq:standard-sequence}
 \boldsymbol\std^\sigma
 =(\std_{\sigma(1)},\ldots,\std_{\sigma(n)}).
\end{equation}
Thus $\sigma^{-1}(v)$ is the position of $\std_v$ in
$\boldsymbol\std^\sigma$. Define the total order $\leq_\sigma$ on
$\{1,\ldots,n\}$ by
\begin{equation}\label{eq:order}
 u\leq_\sigma v\quad\Longleftrightarrow\quad
 \sigma^{-1}(u)\leq\sigma^{-1}(v).
\end{equation}
Write $u<_\sigma v$ if $u\leq_\sigma v$ and $u\neq v$.
The least and greatest weights are $\sigma(1)$ and $\sigma(n)$,
respectively. Thus smaller weights occur earlier in
\eqref{eq:standard-sequence}. By \eqref{eq:exceptional-sequence},
the sequence $\boldsymbol\std^\sigma$ is exceptional precisely when
\begin{equation}\label{eq:semiorthogonal}
 \Hom_{\D_{\fd}(A)}(\std_w,\std_v[t])=0
 \qquad(v<_\sigma w,\ t\in\ZZ).
\end{equation}
The sequence $\boldsymbol\std^\sigma$ is \defn{full in $\D_{\fd}(A)$}
if its terms thickly generate $\D_{\fd}(A)$.

\begin{definition}[$m$-qh dga]
\label{def:shifted}
An algebra $A$ satisfying \eqref{eq:standing} is an
\defn{$m$-qh dga} with
respect to $\sigma$ if there are objects $\std_v\in\H_A^m$ such that:
\begin{enumerate}[label=\textup{(Q\arabic*)},ref=Q\arabic*]
 \item\label{ax:exceptional}
 $\boldsymbol\std^\sigma$ is an exceptional sequence;
 \item\label{ax:filtration}
 for every vertex $v$, there is a conflation
 \begin{equation}\label{eq:presentation}
 K_v\longrightarrow P_v\xrightarrow{p_v}\std_v\dashrightarrow
 \qquad\text{in }\H_A^m,
 \end{equation}
 with
 \begin{equation}\label{eq:kernel}
 \begin{aligned}
 K_v&\in\Filt_{\H_A^m}(\mathcal S_{>v}^{m}),\\
 \mathcal S_{>v}^{m}
 &=\{\std_w[s]:v<_\sigma w,\ 0\leq s<m,
                         \ \std_w[s]\in\H_A^m\}.
 \end{aligned}
 \end{equation}
\end{enumerate}
\end{definition}

The shifts in \eqref{eq:kernel} are cohomological. The membership
$\std_w[s]\in\H_A^m$ is required for each factor, not only for $K_v$.
For $v=\sigma(n)$, the set $\mathcal S_{>v}^{m}$ is empty, so
$K_v=0$ and $p_v:P_v\xrightarrow{\sim}\std_v$.
A particular filtration \eqref{eq:filtration} is a witness for
\eqref{eq:kernel}, rather than additional chosen structure.

\section{Consequences of the axioms}\label{sec:consequences}

Throughout this section, fix a structure satisfying
Definition~\ref{def:shifted}.

\begin{proposition}[Perfection and fullness]\label{prop:fullness}
Conditions~\ref{ax:exceptional} and \ref{ax:filtration} imply
\begin{equation}\label{eq:fullness}
 \std_v\in\per A\quad(1\leq v\leq n),\qquad
 \thick_{\D_{\fd}(A)}(\std_1,\ldots,\std_n)
 =\per A=\D_{\fd}(A).
\end{equation}
\end{proposition}
\begin{proof}
\noindent\textit{Step 1: generation of $\per A$.}
Starting with $\std_{\sigma(n)}\cong P_{\sigma(n)}$, descending
induction in $\leq_\sigma$ and \eqref{eq:kernel} show that $K_v$ and
$\std_v$ are perfect. Conversely, \eqref{eq:presentation} and
\eqref{eq:kernel} put every $P_v$ in
$\thick_{\D_{\fd}(A)}(\std_1,\ldots,\std_n)$.

\noindent\textit{Step 2: generation of $\D_{\fd}(A)$.}
For $X\in\D_{\fd}(A)$, put $X_n=X$. For $a=n,\ldots,1$, take the
evaluation triangle
\begin{equation}\label{eq:evaluation}
 \begin{gathered}
 V_a=\bigoplus_{t\in\ZZ}
 \Hom_{\D_{\fd}(A)}(\std_{\sigma(a)},X_a[t])[-t],\\
 V_a\otimes_{\Bbbk}\std_{\sigma(a)}
 \xrightarrow{\mathrm{ev}_a}X_a\longrightarrow X_{a-1}
 \longrightarrow (V_a\otimes_{\Bbbk}\std_{\sigma(a)})[1].
 \end{gathered}
\end{equation}
Each $V_a$ is a complex with zero differential.
The sums are finite: the first argument is perfect, and the second
has finite-dimensional total cohomology. Exceptionality gives
$\Hom_{\D_{\fd}(A)}(\std_{\sigma(a)},X_{a-1}[t])=0$ for all $t$;
\eqref{eq:semiorthogonal} preserves the vanishings for positions
$b>a$. Thus $\Hom_{\D_{\fd}(A)}(P_v,X_0[t])=0$ for every $v,t$.
Since $\Hom_{\D_{\fd}(A)}(P_v,X_0[t])=H^t(X_0)e_v$, we obtain
$X_0=0$. The triangles \eqref{eq:evaluation} now express $X$ in
$\thick_{\D_{\fd}(A)}(\std_1,\ldots,\std_n)$.
\end{proof}

\begin{corollary}[Smoothness]\label{cor:smooth}
Every algebra satisfying Definition~\ref{def:shifted} is
\defn{homologically smooth}: its diagonal bimodule belongs to
$\per(A^{\op}\otimes_{\Bbbk}A)$.
\end{corollary}
\begin{proof}
Proposition~\ref{prop:fullness} makes every simple $B$-module $S_v$
perfect over $A$. Since $B$ is split basic,
$B/\operatorname{rad}B\cong\Bbbk^n$ is separable over $\Bbbk$.
The smoothness criterion \cite[Theorem 5.8]{RS22} applies.
\end{proof}

The \defn{left-dual sequence} of $\boldsymbol\std^\sigma$ is the
exceptional sequence
$(\costd_{\sigma(n)},\ldots,\costd_{\sigma(1)})$, listed in decreasing
order for $\leq_\sigma$, with terms in
$\thick_{\D_{\fd}(A)}(\std_1,\ldots,\std_n)$, satisfying
\begin{equation}\label{eq:pairing}
 \Hom_{\D_{\fd}(A)}(\std_u,\costd_v[t])\cong
 \begin{cases}\Bbbk,&u=v,\ t=0,\\0,&\text{otherwise}
 \end{cases}
 \qquad(1\leq u,v\leq n,\ t\in\ZZ).
\end{equation}
This is the convention of \cite[\S2.3]{BB26}.

\begin{proposition}[Left-dual objects]\label{prop:costandards}
The sequence $\boldsymbol\std^\sigma$ has a unique left-dual sequence
up to isomorphism of its terms. Its objects satisfy
\begin{equation}\label{eq:costandard-window}
 \costd_v\in\H_A^m[1-m],\qquad
 \H_A^m[1-m]=\D_{\fd}(A)^{[0,m-1]}.
\end{equation}
\end{proposition}
\begin{proof}
Perfection and properness make all mutation triangles finite.
For $v=\sigma(a)$, successively take evaluation cones of
$\std_v$ against
$\std_{\sigma(a-1)},\ldots,\std_{\sigma(1)}$; for $a=1$ take
$\costd_{\sigma(1)}=\std_{\sigma(1)}$.
Semiorthogonality gives \eqref{eq:pairing}, and the mutation
construction gives the reversed exceptional sequence
$(\costd_{\sigma(n)},\ldots,\costd_{\sigma(1)})$
\cite[\S2.3]{BB26}.
Proposition~\ref{prop:fullness} and the characterisation by
\eqref{eq:pairing} for a full exceptional sequence give uniqueness
\cite[\S2.3]{BB26}.

For every factor $\std_w[s]\in\mathcal S_{>v}^{m}$,
\begin{equation}\label{eq:factor-pairing}
 \Hom_{\D_{\fd}(A)}(\std_w[s],\costd_u[t])\cong
 \begin{cases}\Bbbk,&w=u,\ t=s,\\0,&\text{otherwise}.
 \end{cases}
\end{equation}
The long exact sequences from the filtrations of $K_v$ therefore give
$\Hom_{\D_{\fd}(A)}(K_v,\costd_u[t])=0$ for
$t\notin[0,m-1]$. Applying $\Hom_{\D_{\fd}(A)}(-,\costd_u[t])$
to \eqref{eq:presentation} gives
\[
 H^t(\costd_u)e_v
 =\Hom_{\D_{\fd}(A)}(P_v,\costd_u[t])=0
 \qquad(t\notin[0,m-1]).
\]
Summing over $v$ proves \eqref{eq:costandard-window}.
\end{proof}

The fullness of $\boldsymbol\std^\sigma$, existence of its left-dual
sequence $(\costd_{\sigma(n)},\ldots,\costd_{\sigma(1)})$, and
\eqref{eq:costandard-window} are consequences of
Definition~\ref{def:shifted}.

\begin{proposition}[Standards determined by the ordered projectives]
\label{prop:canonical}
For $1\leq v\leq n$, put
\begin{equation}\label{eq:tail}
 \mathcal Q_v=\thick_{\D_{\fd}(A)}\{P_w:v<_\sigma w\},\qquad
 e_{>v}^\sigma=\sum_{v<_\sigma w}e_w.
\end{equation}
Write $\mathcal Q_v^\perp$ for the objects $X\in\D_{\fd}(A)$ such that
$\Hom_{\D_{\fd}(A)}(Y,X[t])=0$ for every $Y\in\mathcal Q_v$ and
$t\in\ZZ$. Then
\begin{equation}\label{eq:localisation}
 \begin{gathered}
 K_v\in\mathcal Q_v,\qquad \std_v\in\mathcal Q_v^\perp,\\
 p_v^*:\Hom_{\D_{\fd}(A)}(\std_v,X[t])
 \xrightarrow{\sim}\Hom_{\D_{\fd}(A)}(P_v,X[t])
 \qquad(X\in\mathcal Q_v^\perp,\ t\in\ZZ).
 \end{gathered}
\end{equation}
Consequently $(\std_v,p_v)$ is determined by $(A,\sigma)$ up to a
unique isomorphism compatible with $p_v$, and $K_v$ is determined up to
isomorphism. Moreover,
\begin{equation}\label{eq:support}
 \begin{gathered}
 H^t(\std_v)e_w=0\quad(v<_\sigma w,\ t\in\ZZ),\\
 H^t(\std_v)e_v\cong
 \begin{cases}\Bbbk,&t=0,\\0,&t\neq0,
 \end{cases}
 \qquad
 H^0(\std_v)\cong e_vB/e_vBe_{>v}^\sigma B.
 \end{gathered}
\end{equation}
The morphism $H^0(p_v):e_vB\to H^0(\std_v)$ is a projective cover
with simple top $S_v$.
\end{proposition}
\begin{proof}
The induction in Step~1 of Proposition~\ref{prop:fullness}, restricted
to weights $w>_\sigma v$, gives
$\mathcal Q_v=\thick_{\D_{\fd}(A)}\{\std_w:v<_\sigma w\}$.
Equations \eqref{eq:kernel} and \eqref{eq:semiorthogonal} imply
$K_v\in\mathcal Q_v$ and $\std_v\in\mathcal Q_v^\perp$.
Applying $\Hom_{\D_{\fd}(A)}(-,X[t])$ to
\eqref{eq:presentation} proves \eqref{eq:localisation} and its
uniqueness assertion. Taking $X=\std_v$ gives
$\Hom_{\D_{\fd}(A)}(P_v,\std_v[t])\cong
\Hom_{\D_{\fd}(A)}(\std_v,\std_v[t])$; the other vertex vanishings
follow from $\std_v\in\mathcal Q_v^\perp$.

For the degree-zero formula, set $D_v^0=e_vB/e_vBe_{>v}^\sigma B$.
The module $D_v^0$ belongs to $\mathcal Q_v^\perp$, so the quotient
$P_v\to D_v^0$ factors through $p_v$ by \eqref{eq:localisation}.
Conversely, the support vanishings and the epimorphism $H^0(p_v)$
give a factorisation $e_vB\to D_v^0\to H^0(\std_v)$.
The two induced maps between $D_v^0$ and $H^0(\std_v)$ are inverse,
as can be checked after precomposition with their epimorphisms from
$e_vB$. Finally, $H^0(\std_v)e_v\cong\Bbbk$, so $H^0(\std_v)$ is a
non-zero quotient of the indecomposable projective $e_vB$.
Thus $H^0(p_v)$ is a projective cover with top $S_v$.
\end{proof}

\begin{corollary}[Classical specialisation]\label{cor:classical}
For $m=1$, Definition~\ref{def:shifted} is equivalent to classical
quasi-heredity of $B=H^0(A)$ with respect to $\leq_\sigma$.
\end{corollary}
\begin{proof}
The standard $t$-structure identifies $\H_A^1$ with $\mod B$, and
$A$ is quasi-isomorphic to $B$ under \eqref{eq:standing} with $m=1$.

\noindent\underline{Shifted axioms $\Rightarrow$ classical axioms.}
Only $s=0$ occurs in \eqref{eq:kernel}. The standard objects are the
usual standard modules by \eqref{eq:support}, and
\eqref{eq:presentation} becomes a short exact sequence whose kernel
is filtered by standards of larger weights. These are the classical
projective-filtration conditions \cite[\S3]{Kra17}.

\noindent\underline{Classical axioms $\Rightarrow$ shifted axioms.}
Classical standard modules satisfy \eqref{eq:presentation} and
\eqref{eq:kernel} with $s=0$. They form an exceptional sequence
\cite[Lemma 2.2]{BB26}, so condition~\ref{ax:exceptional} holds.
\end{proof}

\begin{example}[The order for two classical standards]\label{ex:classical-order}
Take $A=B=\Bbbk(1\xrightarrow{a}2)$ concentrated in degree zero,
with $e_1a=a=ae_2$, $m=1$, and $\sigma=\mathrm{id}$. Thus $1<_\sigma2$.
Let $S_v$ be the simple top of $P_v=e_vB$, for $v=1,2$. Then
\begin{equation}\label{eq:classical-order-objects}
 \std_1=S_1,\qquad \std_2=S_2=P_2,\qquad
 \costd_1=S_1,\qquad \costd_2=P_1.
\end{equation}
The projective presentation
\begin{equation}\label{eq:classical-order-extension}
 0\longrightarrow S_2\longrightarrow P_1\longrightarrow S_1
 \longrightarrow0
\end{equation}
gives
\begin{equation}\label{eq:classical-order-hom}
 \Hom_{\D_{\fd}(B)}(S_2,S_1[t])=0\quad(t\in\ZZ),\qquad
 \Hom_{\D_{\fd}(B)}(S_1,S_2[1])\cong\Bbbk.
\end{equation}
Consequently $(\std_1,\std_2)=(S_1,S_2)$ is full exceptional,
whereas $(\std_2,\std_1)$ is not exceptional. The left-dual sequence
of $(\std_1,\std_2)$ specified by \eqref{eq:pairing} is
$(\costd_2,\costd_1)=(P_1,S_1)$.
\end{example}

\section{Graded path dg algebras}\label{sec:paths}

Let $Q$ be a finite acyclic quiver on $\{1,\ldots,n\}$, with all arrow
degrees at most zero. Paths are concatenated from left to right:
$e_ua=a=ae_v$ for an arrow $a:u\to v$. Let
$A=(\Bbbk Q,d_A)$, where $d_A$ is any graded derivation of degree $1$
with $d_A^2=0$. All vertex idempotents are closed. For a path $p$, let
$|p|$ denote its degree, including $|e_v|=0$, and put
\begin{equation}\label{eq:path-width}
 M=1-\min\{|p|:p\text{ is a path in }Q\}.
\end{equation}
Thus the underlying graded algebra is concentrated in degrees
$[1-M,0]$.

\begin{theorem}[Every ordering]\label{thm:paths}
For every $\sigma\in S_n$, define
\begin{equation}\label{eq:path-objects}
 B_v=A/Ae_{>v}^\sigma A,\qquad
 \std_v=e_vB_v,\qquad K_v=e_vAe_{>v}^\sigma A.
\end{equation}
The objects \eqref{eq:path-objects} give an $M$-qh dga.
The same presentations and filtrations give an $m$-qh dga
for every $m\geq M$.
\end{theorem}
\begin{proof}
\noindent\textit{Step 1: the kernel filtration.}
Fix $v=\sigma(i)$. If $i=n$, then $K_v=0$ and its filtration is empty.
For $i<n$, define dg ideals and dg submodules by
\begin{equation}\label{eq:path-tower}
 \begin{gathered}
 J_j=A\left(\sum_{b=j}^n e_{\sigma(b)}\right)A,\quad
 F_j=e_vJ_j,\quad J_{n+1}=0,\\
 0=F_{n+1}\subseteq F_n\subseteq\cdots\subseteq F_{i+1}=K_v.
 \end{gathered}
\end{equation}
For $i<j\leq n$ and $w=\sigma(j)$, put $L_{vw}=e_vB_we_w$,
a finite complex of vector spaces. Multiplication gives a dg
isomorphism
\begin{equation}\label{eq:path-layer}
 L_{vw}\otimes_{\Bbbk}\std_w
 \xrightarrow{\sim}F_j/F_{j+1},\qquad p\otimes q\longmapsto pq.
\end{equation}
Indeed, a basis path in $F_j/F_{j+1}$ has greatest vertex $w$ in
$\leq_\sigma$. Acyclicity permits a unique cut at the sole occurrence
of $w$. The Leibniz rule makes \eqref{eq:path-layer} a morphism of
complexes. Splitting $L_{vw}$ over $\Bbbk$ gives
\begin{equation}\label{eq:path-factors}
 F_j/F_{j+1}\simeq
 \bigoplus_{h=1-M}^{0}
 \std_w[-h]^{\oplus\dim_{\Bbbk}H^h(L_{vw})}
 \qquad\text{in }\D_{\fd}(A).
\end{equation}

\noindent\textit{Step 2: every factor stays in $\H_A^M$.}
If $H^h(L_{vw})\neq0$, there is a prefix path $p:v\to w$ of degree
$h$. For every basis path $q$ of $\std_w$, the concatenation $pq$ is
a path, and $h+|q|\geq1-M$. Thus
$\std_w[-h]\in\H_A^M$ and $0\leq-h<M$.
All submodules and quotients in \eqref{eq:path-tower} are themselves
concentrated in $[1-M,0]$. Their short exact sequences induce
conflations in $\H_A^M$; refining \eqref{eq:path-factors} gives
\eqref{eq:kernel}. The short exact sequence
$0\to K_v\to P_v\to\std_v\to0$ gives \eqref{eq:presentation}.

\noindent\textit{Step 3: exceptionality.}
Fix $v$. The path basis gives $\std_ve_w=0$ for $v<_\sigma w$ and
$\std_ve_v=\Bbbk$ in degree zero. Descending induction on
$w>_\sigma v$, using the presentations of $\std_w$ and the
filtrations of $K_w$, gives
$\Hom_{\D_{\fd}(A)}(\std_w,\std_v[t])=0$ for every $t$.
The filtration of $K_v$ then gives
$\Hom_{\D_{\fd}(A)}(K_v,\std_v[t])=0$. Consequently
\[
 \Hom_{\D_{\fd}(A)}(\std_v,\std_v[t])
 \cong\Hom_{\D_{\fd}(A)}(P_v,\std_v[t])
 =H^t(\std_ve_v),
\]
which is $\Bbbk$ for $t=0$ and zero otherwise.
For $m\geq M$, every object in the presentations and filtrations still
belongs to $\H_A^m$, and every factor shift still satisfies $0\leq s<m$.
\end{proof}

\begin{example}[A graded quiver of type $A_4$]\label{ex:a4}
Take $Q:1\xrightarrow{a}2\xrightarrow{b}3\xrightarrow{c}4$, with
$|a|=|b|=|c|=-1$ and $d_A=0$. Then $M=4$.
Choose $(\sigma(1),\sigma(2),\sigma(3),\sigma(4))=(2,4,1,3)$,
and put $U=P_1/(abA)$. The standard objects and their kernels are
\begin{equation}\label{eq:a4}
 \begin{array}{c|c|c}
 v&\std_v&K_v\\\hline
 1&U&abA\cong\std_3[2]\\
 2&S_2&bA\cong\std_3[1]\\
 3&P_3&0\\
 4&S_4&0
 \end{array}
 \qquad
 \boldsymbol\std^\sigma=(S_2,S_4,U,P_3).
\end{equation}
The object $P_3$ has cohomology $S_3$ in degree $0$ and $S_4$ in
degree $-1$. Thus the two non-zero kernels in \eqref{eq:a4} occupy
degrees $[-3,-2]$ and $[-2,-1]$, respectively, and both are single
allowed factors in $\H_A^4$.
\end{example}

\begin{example}[A graded quiver of type $D_4$]\label{ex:d4}
Take $Q:1\xrightarrow{a}2\xrightarrow{b}3$, $2\xrightarrow{c}4$,
with $|a|=|c|=-1$, $|b|=0$, and $d_A=0$. Then $M=3$.
Choose $(\sigma(1),\sigma(2),\sigma(3),\sigma(4))=(1,3,2,4)$,
and put $V=P_2/cA$. Then
\begin{equation}\label{eq:d4}
 \begin{array}{c|c|c}
 v&\std_v&K_v\\\hline
 1&S_1&aA\cong P_2[1]\\
 2&V&cA\cong\std_4[1]\\
 3&S_3&0\\
 4&S_4&0
 \end{array}
 \qquad
 \boldsymbol\std^\sigma=(S_1,S_3,V,S_4).
\end{equation}
The kernel $K_1$ has the two-factor filtration
\begin{equation}\label{eq:d4-filtration}
 0\subset acA\subset aA,\qquad
 acA\cong\std_4[2],\qquad aA/acA\cong\std_2[1].
\end{equation}
Both factors belong to $\H_A^3$, with $1<_\sigma4$ and
$1<_\sigma2$. The filtration \eqref{eq:d4-filtration} refines the
shifted-projective description of $K_1$ into shifted standard factors.
\end{example}

\begin{example}[A differential and the required width]\label{ex:differential}
Take arrows $a:1\to2$, $b:2\to3$, and $c:1\to3$, with
\begin{equation}\label{eq:differential}
 |a|=|b|=0,\qquad |c|=-1,\qquad
 d_A(c)=ab,\qquad d_A(a)=d_A(b)=0.
\end{equation}
Here $M=2$, but $H^q(A)=0$ for $q\neq0$ and
$B=H^0(A)=\Bbbk(1\xrightarrow{a}2\xrightarrow{b}3)/(ab)$.
For $(\sigma(1),\sigma(2),\sigma(3))=(1,3,2)$,
\begin{equation}\label{eq:differential-standards}
 (\std_1,\std_2,\std_3)=(S_1,P_2,S_3),\qquad
 K_1=\Bbbk\{a,ab,c\},\qquad K_2=K_3=0.
\end{equation}
There is a filtration by dg submodules
\begin{equation}\label{eq:differential-filtration}
 0\subset aA\subset K_1,\qquad
 aA\cong\std_2,\qquad K_1/aA\cong\std_3[1].
\end{equation}
Theorem~\ref{thm:paths} gives a structure of width $2$.
Although $K_1\simeq S_2$ and all the objects in
\eqref{eq:differential-standards} lie in $\H_A^1$, the factor
$\std_3[1]$ does not. There is no width-$1$ structure for this
ordering: Proposition~\ref{prop:canonical} fixes the standard objects.
Write $[X]$ for the class of a module $X$ in $K_0(\mod B)$, whose
basis is $[S_1],[S_2],[S_3]$. A filtration of $S_2$ by $P_2$ and $S_3$
in $\mod B$ would require
\[
 [S_2]=x[P_2]+y[S_3]=x[S_2]+(x+y)[S_3],
 \qquad x,y\in\ZZ_{\geq0},
\]
which forces $x=1$ and $y=-1$, a contradiction.
\end{example}

\section{Filtrations, width, and opposite algebras}\label{sec:filtrations}

Fix a structure satisfying Definition~\ref{def:shifted}.
For an integer $\ell\geq1$, define
\begin{equation}\label{eq:allowed}
 I_w^\ell=\{s\in\ZZ:0\leq s<\ell,\ \std_w[s]\in\H_A^\ell\},
 \qquad
 \mathcal S_{>v}^\ell=\{\std_w[s]:v<_\sigma w,\ s\in I_w^\ell\}.
\end{equation}
For $1\leq v,w\leq n$ and $t\in\ZZ$, put
\begin{equation}\label{eq:numerical}
 b_{vw,t}=\dim_{\Bbbk}\Hom_{\D_{\fd}(A)}(P_v,\costd_w[t])
         =\dim_{\Bbbk}H^t(\costd_w)e_v.
\end{equation}

\begin{proposition}[A criterion for shifted filtrations]\label{prop:criterion}
For $X\in\mathcal Q_v$, where $\mathcal Q_v$ is \eqref{eq:tail},
\begin{equation}\label{eq:criterion}
 \begin{split}
 X\in\Filt_{\H_A^\ell}(\mathcal S_{>v}^\ell)
 \quad\Longleftrightarrow\quad&
 \Hom_{\D_{\fd}(A)}(X,\costd_w[t])=0\\[-2pt]
 &\text{for }v<_\sigma w,\ t\notin I_w^\ell.
 \end{split}
\end{equation}
When \eqref{eq:criterion} holds, $X$ has a filtration by decreasing
weights whose weight-$w$ factor is
\begin{equation}\label{eq:weight-block}
 G_w(X)\cong
 \bigoplus_{s\in I_w^\ell}
 \std_w[s]^{\oplus\dim_{\Bbbk}
                    \Hom_{\D_{\fd}(A)}(X,\costd_w[s])}.
\end{equation}
The category $\Filt_{\H_A^\ell}(\mathcal S_{>v}^\ell)$ is closed under
direct summands. For $X=K_v$, the multiplicities in
\eqref{eq:weight-block} are $b_{vw,s}$.
\end{proposition}
\begin{proof}
We prove necessity, construct the filtration for sufficiency, and
then deduce the assertions about multiplicities and direct summands.

\noindent\underline{Necessity.}
By \eqref{eq:factor-pairing}, an allowed factor $\std_u[s]$ has a
non-zero morphism to $\costd_w[t]$ only when $u=w$ and $t=s$.
The long exact sequences of a filtration give the vanishings in
\eqref{eq:criterion}.

\noindent\underline{Sufficiency.}
Put $a=\sigma^{-1}(v)$ and apply the evaluation construction
\eqref{eq:evaluation} to $X$ in the exceptional subcategory
$\mathcal Q_v$. The resulting triangles, combined by the octahedral
axiom, give an ordered tower with successive factors in
$\thick_{\D_{\fd}(A)}(\std_{\sigma(b)})$
for $b=n,n-1,\ldots,a+1$. If $a=n$, then $\mathcal Q_v=0$ and the
tower is empty.
For each weight $w$, exceptionality identifies
$\thick_{\D_{\fd}(A)}(\std_w)$ with $\D^b(\mod\Bbbk)$, so its factor
is a finite direct sum of shifts of $\std_w$.
Applying $\Hom_{\D_{\fd}(A)}(-,\costd_w[t])$ to the tower kills every
factor of weight different from $w$. The pairing \eqref{eq:pairing}
therefore identifies the multiplicity of $\std_w[t]$ with
$\dim_{\Bbbk}\Hom_{\D_{\fd}(A)}(X,\costd_w[t])$.
The vanishings in \eqref{eq:criterion} give precisely
\eqref{eq:weight-block}. Every factor and, by extension closure, every
intermediate object belongs to $\H_A^\ell$.
Splitting the finite direct sums refines the tower to a filtration
by $\mathcal S_{>v}^\ell$.

Both $\mathcal Q_v$ and the vanishing conditions in
\eqref{eq:criterion} are closed under direct summands.
For $v<_\sigma w$, applying
$\Hom_{\D_{\fd}(A)}(-,\costd_w[t])$ to \eqref{eq:presentation} gives
$\Hom_{\D_{\fd}(A)}(K_v,\costd_w[t])
 \cong\Hom_{\D_{\fd}(A)}(P_v,\costd_w[t])$.
This proves the final assertion using \eqref{eq:numerical}.
\end{proof}

Arbitrary filtrations still need not have invariant multiplicities.
For $v<_\sigma w$ and $0\leq s<m-1$ with
$\std_w[s],\std_w[s+1]\in\H_A^m$, the conflations
\begin{equation}\label{eq:cancellation}
 0\longrightarrow\std_w[s]\xrightarrow{1}\std_w[s]\dashrightarrow,
 \qquad
 \std_w[s]\longrightarrow0\longrightarrow\std_w[s+1]\dashrightarrow
\end{equation}
give a two-factor filtration of zero.
In contrast, the multiplicities in the ordered factors
\eqref{eq:weight-block} are fixed by the left-dual pairing.
If $c_{vw,s}$ counts factors $\std_w[s]$ in any filtration of $K_v$,
Euler additivity gives
\begin{equation}\label{eq:signed-multiplicity}
 \sum_{s=0}^{m-1}(-1)^s c_{vw,s}
 =\sum_{s=0}^{m-1}(-1)^s b_{vw,s}
 \qquad(v<_\sigma w).
\end{equation}
Proposition~\ref{prop:criterion} supplies a filtration with
$c_{vw,s}=b_{vw,s}$ for all $v<_\sigma w$ and $0\leq s<m$.

\begin{corollary}[Least width]\label{cor:width}
Suppose $(A,\sigma)$ admits a structure of
Definition~\ref{def:shifted} for at least one width. For its canonical
standard objects and specified left-dual objects, put
\begin{equation}\label{eq:depths}
 \begin{aligned}
 \alpha_A&=-\min\{q:H^q(A)\neq0\},\\
 \alpha_w&=-\min\{q:H^q(\std_w)\neq0\},&
 \beta_w&=\max\{q:H^q(\costd_w)\neq0\}.
 \end{aligned}
\end{equation}
Then $\alpha_A,\alpha_w,\beta_w\geq0$, and the least width is
\begin{equation}\label{eq:minimal-width}
 m_{\min}(A,\sigma)
 =1+\max\left\{\alpha_A,\ \max_{1\leq w\leq n}
                         (\alpha_w+\beta_w)\right\}.
\end{equation}
\end{corollary}
\begin{proof}
Proposition~\ref{prop:canonical} makes $\std_w$ and $\costd_w$
independent of the chosen width. Its proof, together with
\eqref{eq:pairing}, also gives
\begin{equation}\label{eq:dual-support}
 H^t(\costd_w)e_v=0\quad(w<_\sigma v),\qquad
 H^t(\costd_w)e_w\cong
 \begin{cases}\Bbbk,&t=0,\\0,&t\neq0.\end{cases}
\end{equation}
Thus $\beta_w\geq0$, and if $\beta_w>0$, then
$b_{vw,\beta_w}\neq0$ for some $v<_\sigma w$.
Since $H^0(\std_w)\neq0$, membership of a shifted standard is
equivalent to
\[
 \std_w[s]\in\H_A^\ell
 \quad\Longleftrightarrow\quad
 0\leq s\leq\ell-1-\alpha_w
 \qquad(s\geq0).
\]
Necessity in \eqref{eq:criterion} therefore forces
$\alpha_w+\beta_w\leq\ell-1$ whenever width $\ell$ is possible;
when $\beta_w=0$, the same inequality follows from
$\std_w\in\H_A^\ell$. Also $\alpha_A\leq\ell-1$ is necessary.

Conversely, these inequalities put every non-zero factor
$\std_w[s]$ of the ordered kernel tower in $\H_A^\ell$, with
$0\leq s<\ell$. Proposition~\ref{prop:criterion} gives
$K_v\in\Filt_{\H_A^\ell}(\mathcal S_{>v}^\ell)$.
The objects $P_v$ and $\std_v$ also lie in $\H_A^\ell$, so their
presentation triangles are conflations. This proves
\eqref{eq:minimal-width}.
\end{proof}

\begin{proposition}[Injective presentations and opposite algebras]
\label{prop:injectives}
Set $\widehat{\costd}_v=\costd_v[m-1]$ and
$I_v=\dual(Ae_v)[m-1]$. There are conflations
\begin{equation}\label{eq:injective}
 \begin{gathered}
 \widehat{\costd}_v\longrightarrow I_v\longrightarrow L_v
 \dashrightarrow\quad\text{in }\H_A^m,\\
 L_v\in\Filt_{\H_A^m}
 \{\widehat{\costd}_w[-s]:v<_\sigma w,\ 0\leq s<m,\
                         \widehat{\costd}_w[-s]\in\H_A^m\}.
 \end{gathered}
\end{equation}
The algebra $A^{\op}$ satisfies Definition~\ref{def:shifted} for the
same $m$ and $\sigma$, with standards $\dual\costd_v$ and specified
left-dual objects $\dual\std_v$.
\end{proposition}
\begin{proof}
Decompose $I_v$ along the full exceptional sequence
$(\costd_{\sigma(n)},\ldots,\costd_{\sigma(1)})$.
The factors of $I_v$ occur in increasing order in $\leq_\sigma$.
The pairing \eqref{eq:pairing} and the identity
\[
 \Hom_{\D_{\fd}(A)}(\std_w,I_v[t])
 \cong\dual\bigl(H^{1-m-t}(\std_w)e_v\bigr)
\]
identify the weight-$w$ factor as
\begin{equation}\label{eq:injective-factors}
 \bigoplus_{s\geq0}
 \costd_w[m-1-s]^{\oplus\dim_{\Bbbk}H^{-s}(\std_w)e_v}.
\end{equation}
By \eqref{eq:support}, weights $w<_\sigma v$ contribute zero, and
the first non-zero factor is $\widehat{\costd}_v$.
Every occurring $s$ satisfies $s\leq\alpha_w$; hence
$\beta_w+s\leq m-1$ by Corollary~\ref{cor:width}.
Thus all factors in \eqref{eq:injective-factors} belong to $\H_A^m$,
giving \eqref{eq:injective}.

The contravariant equivalence
$\dual(-)[m-1]:\H_A^m\to\H_{A^{\op}}^m$
takes \eqref{eq:injective} to the required projective presentation:
\[
 \dual L_v[m-1]\longrightarrow Ae_v
 \longrightarrow\dual\costd_v\dashrightarrow.
\]
Its kernel factors are $(\dual\costd_w)[s]$ with $v<_\sigma w$.
Vector-space duality reverses the costandard exceptional sequence,
giving $(\dual\costd_{\sigma(1)},\ldots,\dual\costd_{\sigma(n)})$.
Finally,
\[
 \Hom_{\D_{\fd}(A^{\op})}(\dual\costd_u,\dual\std_v[t])
 \cong\Hom_{\D_{\fd}(A)}(\std_v,\costd_u[t])
\]
proves the asserted left-dual pairing.
\end{proof}

\section{Framework for replacing (Q1)}\label{sec:comparison-framework}

In Sections~\ref{sec:proposal-hom}--\ref{sec:proposal-negative}, retain
the standing assumptions \eqref{eq:standing}, the order \eqref{eq:order},
and condition~\ref{ax:filtration}. The objects $\std_v\in\H_A^m$
are not assumed exceptional. Each proposal retains the
\defn{brick condition}
\begin{equation}\label{eq:trial-brick}
 \End_{\H_A^m}(\std_v)\cong\Bbbk\qquad(1\leq v\leq n),
\end{equation}
unless \eqref{eq:trial-brick} is explicitly proved from stronger
conditions. The purpose of replacing condition~\ref{ax:exceptional}
is to obtain homological smoothness as a theorem. Smoothness as an
additional assumption is considered only for comparison in
Section~\ref{sec:proposal-assumed-smooth}.

For a finite-dimensional $B$-module $M$, write $[M:S_w]$ for its
composition multiplicity. Since $B$ is split basic,
\begin{equation}\label{eq:trial-projective-test}
 [M:S_w]=\dim_{\Bbbk}Me_w,\qquad
 \Hom_{\D_{\fd}(A)}(P_w,X[t])\cong H^t(X)e_w
 \quad(X\in\D_{\fd}(A),\ t\in\ZZ).
\end{equation}
The equality
$\Hom_{\H_A^m}(X,Y)=\Hom_{\D_{\fd}(A)}(X,Y)$
applies only when $X,Y\in\H_A^m$.

\begin{proposition}[Generation from the filtration axiom]
\label{prop:trial-generation}
Condition~\ref{ax:filtration} alone implies
\begin{equation}\label{eq:trial-generation}
 \std_v\in\per A\quad(1\leq v\leq n),\qquad
 \thick_{\D_{\fd}(A)}(\std_1,\ldots,\std_n)=\per A.
\end{equation}
\end{proposition}
\begin{proof}
The greatest weight has $K_{\sigma(n)}=0$ and
$\std_{\sigma(n)}\cong P_{\sigma(n)}$. Descending induction using
\eqref{eq:kernel} and \eqref{eq:presentation} makes every $K_v$ and
$\std_v$ perfect. Conversely, the same triangles put every $P_v$ in
the thick subcategory generated by the standards.
\end{proof}

\begin{proposition}[The classical projective test]
\label{prop:trial-classical}
For $m=1$, the brick condition \eqref{eq:trial-brick},
condition~\ref{ax:filtration}, and either of the conditions
\begin{align}
 \Hom_B(P_w,\std_v)&=0\quad(v<_\sigma w),
 \label{eq:trial-classical-projectives}\\
 \Hom_B(\std_w,\std_v)&=0\quad(v<_\sigma w)
 \label{eq:trial-classical-standards}
\end{align}
are equivalent to classical quasi-heredity. Here the equivalence
$\H_A^1\simeq\mod B$ identifies $P_w$ with $e_wB$.
\end{proposition}
\begin{proof}
First prove the equivalence of
\eqref{eq:trial-classical-projectives} and
\eqref{eq:trial-classical-standards}.

\noindent\underline{Projective test $\Rightarrow$ standard test.}
Precomposition with the epimorphism $P_w\twoheadrightarrow\std_w$
embeds $\Hom_B(\std_w,\std_v)$ in $\Hom_B(P_w,\std_v)$.

\noindent\underline{Standard test $\Rightarrow$ projective test.}
The filtration of $K_w$ and
\eqref{eq:trial-classical-standards} give
$\Hom_B(K_w,\std_v)=0$ for $v<_\sigma w$.
Apply $\Hom_B(-,\std_v)$ to
$0\to K_w\to P_w\to\std_w\to0$.

Under either test, $\std_v$ is supported on weights at most $v$,
and $\Hom_B(K_v,\std_v)=0$. Consequently
\[
 \Bbbk\cong\End_B(\std_v)
 \cong\Hom_B(P_v,\std_v)\cong\std_ve_v.
\]
The quotient $P_v\twoheadrightarrow\std_v$ has simple top $S_v$,
with $S_v$ occurring once. The projective kernels are filtered by
standards of larger weights. These are the classical axioms
\cite[\S2.1]{BB26}; classical standards conversely satisfy both tests.
\end{proof}

The brick condition cannot be omitted from
Proposition~\ref{prop:trial-classical}. The ordinary algebra
$B=\Bbbk[x]/(x^2)$ with one standard $\std_1=P_1=B$ and $K_1=0$
satisfies both off-diagonal tests, whose index sets are empty,
but is not quasi-hereditary.

\begin{proposition}[Cohomological support]
\label{prop:trial-support}
Assume condition~\ref{ax:filtration} and
\begin{equation}\label{eq:trial-support}
 H^t(\std_v)e_w=0\qquad(v<_\sigma w,\ t\in\ZZ).
\end{equation}
Then
\begin{align}
 \Hom_{\D_{\fd}(A)}(\std_w,\std_v[t])&=0
       &&(v<_\sigma w,\ t\in\ZZ),\label{eq:trial-backward}\\
 \Hom_{\D_{\fd}(A)}(\std_v,\std_v[t])&\cong H^t(\std_v)e_v
       &&(1\leq v\leq n,\ t\in\ZZ).
       \label{eq:trial-diagonal}
\end{align}
In particular, \eqref{eq:trial-brick} is equivalent to
$[H^0(\std_v):S_v]=1$ for every $v$ under these hypotheses.
\end{proposition}
\begin{proof}
Fix $v$ and descend through the weights $w>_\sigma v$.
Equation~\eqref{eq:trial-projective-test} gives
\[
 \Hom_{\D_{\fd}(A)}(P_w,\std_v[t])=0\qquad(t\in\ZZ).
\]
Once the standards of weight greater than $w$ are orthogonal to
$\std_v$ in every degree, their shifted filtration of $K_w$ gives
\[
 \Hom_{\D_{\fd}(A)}(K_w,\std_v[t])=0\qquad(t\in\ZZ).
\]
The presentation of $\std_w$ proves \eqref{eq:trial-backward}.
The filtration of $K_v$ now gives the same vanishing for $K_v$;
precomposition with $p_v:P_v\to\std_v$ proves
\eqref{eq:trial-diagonal}.
\end{proof}

\begin{example}[A non-smooth dga with one brick standard]
\label{ex:trial-dualnumbers}
Let
\begin{equation}\label{eq:trial-dualnumbers}
 R=\Bbbk[\varepsilon]/(\varepsilon^2),\qquad
 |\varepsilon|=-1,\qquad d_R=0,
 \qquad m=2,\quad \std_1=P_1=R,\quad K_1=0.
\end{equation}
The object $R\in\H_R^2$ is a brick, and
\begin{equation}\label{eq:trial-dualnumbers-self}
 \Hom_{\D_{\fd}(R)}(R,R[t])\cong
 \begin{cases}\Bbbk,&t=0,-1,\\0,&\text{otherwise}.
 \end{cases}
\end{equation}
To verify non-smoothness, put $S=R/(\varepsilon)$.
A semifree resolution $F\to S$ is given by
\begin{equation}\label{eq:trial-dualnumbers-resolution}
 F=\bigoplus_{a\geq0}x_aR,\quad |x_a|=-2a,\quad
 d_F(x_0)=0,\quad d_F(x_a)=x_{a-1}\varepsilon\ (a\geq1),
\end{equation}
where the augmentation sends $x_0$ to $1$ and $x_a$ to $0$ for $a>0$.
Each $x_a\varepsilon$ is the boundary of $x_{a+1}$, so $H^*(F)=S$.
Applying the dg Hom complex into $S$ gives
\begin{equation}\label{eq:trial-dualnumbers-ext}
 \Hom_{\D_{\fd}(R)}(S,S[t])\cong
 \begin{cases}\Bbbk,&t\in2\ZZ_{\geq0},\\0,&\text{otherwise}.
 \end{cases}
\end{equation}
Thus $S\notin\per R$, since a perfect object has finite-dimensional
total derived Hom into $S$. Smoothness would put every
cohomologically finite-dimensional module in $\per R$: tensor a
finite construction of the diagonal bimodule with that module.
Therefore $R$ is not smooth and $\per R\subsetneq\D_{\fd}(R)$.
\end{example}

\section{Bricks and backward Hom-orthogonality}
\label{sec:proposal-hom}

Replace condition~\ref{ax:exceptional} by
\eqref{eq:trial-brick} and
\begin{equation*}\tag{H}\label{eq:proposal-hom}
 \Hom_{\H_A^m}(\std_w,\std_v)=0
 \qquad(v<_\sigma w).
\end{equation*}
Together with condition~\ref{ax:filtration},
\eqref{eq:proposal-hom} has the correct classical specialisation by
Proposition~\ref{prop:trial-classical}. For $m>1$, a factor
$\std_u[s]$ of $K_w$ contributes
$\Hom_{\D_{\fd}(A)}(\std_u,\std_v[-s])$ to
$\Hom_{\D_{\fd}(A)}(K_w,\std_v)$.
Equation~\eqref{eq:proposal-hom} controls only $s=0$.

\begin{example}[Backward extensions on a graded fork]
\label{ex:trial-fork}
Take
\[
 A=\Bbbk(1\xrightarrow{a}3\xleftarrow{b}2),\qquad
 |a|=|b|=-1,\quad d_A=0,\quad m=2,
\]
with $e_1a=a=ae_3$, $e_2b=b=be_3$, and order $1<2<3$.
Choose
\begin{equation}\label{eq:trial-fork-standards}
 (\std_1,\std_2,\std_3)=(P_1,S_2,P_3),\qquad
 (K_1,K_2,K_3)=(0,P_3[1],0).
\end{equation}
The short exact sequence of dg modules
\begin{equation}\label{eq:trial-fork-presentation}
 0\longrightarrow P_3[1]\xrightarrow{b}P_2
 \longrightarrow S_2\longrightarrow0
\end{equation}
induces the required conflation in $\H_A^2$.
Every $\std_v$ is individually exceptional and
$H^0(\std_v)=S_v$. Nevertheless,
\begin{equation}\label{eq:trial-fork-homs}
 \Hom_{\D_{\fd}(A)}(\std_w,\std_v[t])\cong
 \begin{cases}
 \Bbbk,&(w,v,t)=(3,1,-1)\text{ or }(2,1,1),\\
 0,&w>v\text{ otherwise}.
 \end{cases}
\end{equation}
Indeed, $\Hom_{\D_{\fd}(A)}(P_2,P_1[t])=0$ for every $t$;
applying derived Hom into $P_1$ to
\eqref{eq:trial-fork-presentation} computes the pair $(2,1)$.
Equation~\eqref{eq:trial-fork-homs} proves
\eqref{eq:proposal-hom} but disproves backward orthogonality in
degrees $1$ and $-1$. The algebra is smooth; the failure concerns
the proposed standards and their order.
\end{example}

\begin{example}[Individual exceptionality does not ensure smoothness]
\label{ex:trial-cycle}
Let $A$ have vertices $1,2$, arrows $a:1\to2$ and $b:2\to1$,
with $|a|=|b|=-1$, $d_A=0$, and relations $ab=ba=0$.
Paths are concatenated from left to right. Take $m=2$, order $1<2$,
and $\std_i=P_i$, $K_i=0$ for $i=1,2$.
Since $e_iAe_i=\Bbbk$, both standards are exceptional, and
\eqref{eq:proposal-hom} holds. But
\[
 \Hom_{\D_{\fd}(A)}(P_2,P_1[-1])=\Bbbk a.
\]
An alternating semifree resolution $F\to S_1$ has generators $x_j$
of degree $-2j$, at vertex $1$ for even $j$ and at vertex $2$ for
odd $j$. Its differential is
$d_F(x_j)=x_{j-1}a$ for odd $j$ and
$d_F(x_j)=x_{j-1}b$ for even $j>0$; $d_F(x_0)=0$.
The augmentation $F\to S_1$ sends $x_0$ to a generator of $S_1$
and every $x_j$, $j>0$, to zero.
The relations $ab=ba=0$ make $d_F^2=0$, and the augmentation is a
quasi-isomorphism. Hence
\begin{equation}\label{eq:trial-cycle-ext}
 \Hom_{\D_{\fd}(A)}(S_1,S_1[4j])\cong\Bbbk
 \qquad(j\geq0).
\end{equation}
The module $S_1$ is not perfect, so $A$ is not smooth.
Thus even supplementing \eqref{eq:proposal-hom} with individual
exceptionality leaves a non-smooth example.
\end{example}

\begin{center}\small
\begin{tabularx}{\linewidth}{@{}>{\raggedright\arraybackslash}p{.36\linewidth}>{\raggedright\arraybackslash}X@{}}
\toprule
Property under \eqref{eq:proposal-hom} & Result with
\eqref{eq:trial-brick} and condition~\ref{ax:filtration} \\\midrule
Classical specialisation & Yes; Proposition~\ref{prop:trial-classical}.\\
Bricks; backward degree-zero Hom & Required.\\
Standards generate $\per A$ & Yes; Proposition~\ref{prop:trial-generation}.\\
Backward Hom in every degree & No; Example~\ref{ex:trial-fork}.\\
Self-exceptionality & Not forced; Example~\ref{ex:trial-dualnumbers}.\\
Standards generate $\D_{\fd}(A)$ & Not forced; Example~\ref{ex:trial-dualnumbers}.\\
Smoothness as a consequence & No, even with exceptional terms;
Example~\ref{ex:trial-cycle}.\\
\bottomrule
\end{tabularx}
\end{center}

\section{Cohomological support with a simple diagonal}
\label{sec:proposal-strong}

Replace condition~\ref{ax:exceptional} by
\begin{equation*}\tag{C}\label{eq:proposal-strong}
 \begin{gathered}
 H^q(\std_v)e_w=0\quad(v<_\sigma w,\ q\in\ZZ),\\
 H^q(\std_v)e_v\cong
 \begin{cases}\Bbbk,&q=0,\\0,&q\neq0\end{cases}
 \quad(1\leq v\leq n).
 \end{gathered}
\end{equation*}
Equivalently, $H^0(\std_v)$ has weights at most $v$, with $S_v$
occurring once, and $H^q(\std_v)$ has weights strictly smaller than
$v$ for $q<0$. No Hom- or Ext-vanishing between standards is imposed.
In the formulation \eqref{eq:proposal-strong} and
condition~\ref{ax:filtration}, neither brickness, exceptionality,
fullness, nor smoothness is an additional axiom.

\begin{proposition}[An equivalent composition-factor definition]
\label{prop:trial-strong}
Under condition~\ref{ax:filtration},
\eqref{eq:proposal-strong} is equivalent to
condition~\ref{ax:exceptional}. In particular, the brick condition,
fullness in $\D_{\fd}(A)=\per A$, and smoothness follow.
\end{proposition}
\begin{proof}
\noindent\underline{\eqref{eq:proposal-strong} $\Rightarrow$
condition~\ref{ax:exceptional}.}
Proposition~\ref{prop:trial-support} proves backward orthogonality
in every degree. Its diagonal identity
\eqref{eq:trial-diagonal} and \eqref{eq:proposal-strong} prove
exceptionality of each $\std_v$.

\noindent\underline{Condition~\ref{ax:exceptional} $\Rightarrow$
\eqref{eq:proposal-strong}.}
Apply Proposition~\ref{prop:canonical}, specifically
\eqref{eq:support}. Proposition~\ref{prop:fullness} and
Corollary~\ref{cor:smooth} give the remaining assertions.
\end{proof}

\begin{proposition}[Smoothness directly from composition factors]
\label{prop:trial-strong-direct-smooth}
Under the standing assumptions \eqref{eq:standing},
\eqref{eq:proposal-strong} and condition~\ref{ax:filtration} imply
\begin{equation}\label{eq:trial-strong-simple-perfect}
 S_v\in\per A\quad(1\leq v\leq n),\qquad
 \D_{\fd}(A)=\per A,\qquad
 A\in\per(A^{\op}\otimes_{\Bbbk}A).
\end{equation}
\end{proposition}
\begin{proof}
\noindent\textit{Step 1: descending induction makes the standards perfect.}
Condition~\ref{ax:filtration} gives
$\std_{\sigma(n)}\cong P_{\sigma(n)}$.
For each $v$, once the standards $\std_w$ with $w>_\sigma v$ are
perfect, their shifted filtration \eqref{eq:kernel} makes $K_v$
perfect. The triangle \eqref{eq:presentation} then makes $\std_v$
perfect.

\noindent\textit{Step 2: ascending induction makes the simples perfect.}
Since $H^1(K_v)=0$, the morphism
$H^0(p_v):e_vB\twoheadrightarrow H^0(\std_v)$ is surjective.
Condition~\eqref{eq:proposal-strong} gives
$[H^0(\std_v):S_v]=1$, so $H^0(\std_v)$ has simple top $S_v$.
Let $q_v:H^0(\std_v)\twoheadrightarrow S_v$ be a top quotient,
let $a_v:\std_v\to S_v$ be the composite of the canonical
truncation morphism $\std_v\to H^0(\std_v)$ with $q_v$, and take
the triangle
\begin{equation}\label{eq:trial-strong-simple-triangle}
 T_v\longrightarrow\std_v\xrightarrow{a_v}S_v
 \longrightarrow T_v[1].
\end{equation}
The cohomology of $T_v$ satisfies
\begin{equation}\label{eq:trial-strong-simple-fibre}
 H^q(T_v)\cong
 \begin{cases}
 H^q(\std_v),&q<0,\\
 \Ker q_v,&q=0,\\
 0,&q>0.
 \end{cases}
\end{equation}
By \eqref{eq:proposal-strong}, every composition factor of every
$H^q(T_v)$ is $S_w$ with $w<_\sigma v$. Finite truncation triangles
and composition series therefore give
$T_v\in\thick_{\D_{\fd}(A)}\{S_w:w<_\sigma v\}$.
For the least weight, $T_{\sigma(1)}=0$ and
$S_{\sigma(1)}\cong\std_{\sigma(1)}$ is perfect. Ascending
induction and \eqref{eq:trial-strong-simple-triangle} make every
$S_v$ perfect. Every object of $\D_{\fd}(A)$ is obtained from the
$S_v$ by finite truncation triangles and composition series, so
$\D_{\fd}(A)=\per A$.

\noindent\textit{Step 3: perfect simples imply smoothness.}
The radical quotient $B/\operatorname{rad}B\cong\Bbbk^n$ is
separable over $\Bbbk$ and perfect over $A$. The criterion
\cite[Theorem 5.8]{RS22} gives
$A\in\per(A^{\op}\otimes_{\Bbbk}A)$.
\end{proof}

The obstruction excluded by \eqref{eq:proposal-strong} is explicit:
for $R=\Bbbk[\varepsilon]/(\varepsilon^2)$, $|\varepsilon|=-1$,
$d_R=0$, and $\std_1=R$, the group
$H^{-1}(\std_1)e_1=\Bbbk\varepsilon$ violates the diagonal condition.
Example~\ref{ex:trial-dualnumbers} proves that removing precisely
that restriction admits a non-smooth algebra. There is no admitted
non-smooth counterexample under \eqref{eq:proposal-strong} and
condition~\ref{ax:filtration}, by Proposition~\ref{prop:trial-strong}.

Negative cohomology of standards is still permitted at smaller
weights. For example, take
$A=\Bbbk(1\xrightarrow{a}2)$ with $|a|=-1$, $d_A=0$, $m=2$,
and order $2<_\sigma1$. Then $\std_1=P_1$, $\std_2=P_2$,
$K_1=K_2=0$, and $H^{-1}(\std_1)=S_2$.

\begin{center}\small
\begin{tabularx}{\linewidth}{@{}>{\raggedright\arraybackslash}p{.36\linewidth}>{\raggedright\arraybackslash}X@{}}
\toprule
Property under \eqref{eq:proposal-strong} & Result with
condition~\ref{ax:filtration} \\\midrule
Classical specialisation & Yes; Corollary~\ref{cor:classical}.\\
Bricks and backward Hom in every degree & Consequences of
Proposition~\ref{prop:trial-support}.\\
Exceptional standard sequence & A consequence, not an axiom.\\
Generation of $\per A$ and $\D_{\fd}(A)$ & Both hold.\\
Smoothness as a consequence & Yes; the direct proof is
Proposition~\ref{prop:trial-strong-direct-smooth}.\\
Negative cohomology of standards & Allowed at strictly smaller weights.\\
Non-smooth algebras excluded & Yes, as a consequence of the axioms.
In particular, the dg dual numbers fail the diagonal requirement
in \eqref{eq:proposal-strong}.\\
Additional restriction compared with (Q1) and (Q2) & None;
Proposition~\ref{prop:trial-strong} proves equivalence.\\
\bottomrule
\end{tabularx}
\end{center}

\section{Allowing the diagonal in negative cohomology}
\label{sec:proposal-weak}

Retain \eqref{eq:trial-brick} and replace
condition~\ref{ax:exceptional} by
\begin{equation*}\tag{C'}\label{eq:proposal-weak}
 H^q(\std_v)e_w=0\qquad(v<_\sigma w,\ q\in\ZZ).
\end{equation*}
Under condition~\ref{ax:filtration},
Proposition~\ref{prop:trial-support} shows that
$[H^0(\std_v):S_v]=1$ follows from brickness. It also proves all
backward derived orthogonality and
\begin{equation}\label{eq:trial-weak-self}
 \Hom_{\D_{\fd}(A)}(\std_v,\std_v[t])
 \cong H^t(\std_v)e_v=0\qquad(t>0).
\end{equation}
Negative self-Hom spaces remain unrestricted by
\eqref{eq:proposal-weak}.

The data
\[
 A=\Bbbk[\varepsilon]/(\varepsilon^2),\quad
 |\varepsilon|=-1,\quad d_A=0,\quad
 m=2,\quad\std_1=P_1=A,\quad K_1=0
\]
satisfy \eqref{eq:trial-brick}, \eqref{eq:proposal-weak}, and
condition~\ref{ax:filtration}. Example~\ref{ex:trial-dualnumbers}
proves that $A$ is not smooth and that
$\operatorname{thick}_{\D_{\fd}(A)}(\std_1)=\per A
\subsetneq\D_{\fd}(A)$.
Even the pairing object cannot be a left-dual exceptional object:
requiring
$\Hom_{\D_{\fd}(A)}(A,N[t])=\Bbbk$ for $t=0$ and zero otherwise
forces $N\simeq\Bbbk$, which is neither perfect nor exceptional by
\eqref{eq:trial-dualnumbers-ext}.

\begin{proposition}[Smoothness and the diagonal]
\label{prop:trial-smooth-diagonal}
Assume \eqref{eq:trial-brick}, \eqref{eq:proposal-weak}, and
condition~\ref{ax:filtration}. Then the following are equivalent:
\begin{enumerate}[label=\textup{(\roman*)}]
 \item $A$ is homologically smooth;
 \item $H^{q}(\std_v)e_v=0$ for all $v$ and all $q<0$;
 \item $\boldsymbol\std^\sigma$ is an exceptional sequence.
\end{enumerate}
\end{proposition}
\begin{proof}
We prove $(\mathrm{i})\Rightarrow(\mathrm{ii})
\Rightarrow(\mathrm{iii})\Rightarrow(\mathrm{i})$.

\noindent\underline{$(\mathrm{i})\Rightarrow(\mathrm{ii})$.}
Put $E_v=\REnd_A(\std_v)$. Equations
\eqref{eq:trial-diagonal} and \eqref{eq:trial-brick} imply that
$E_v$ is proper and connective, with $H^0(E_v)\cong\Bbbk$.
Choose homotopically projective representatives of the standards.
In their dg endomorphism category, the backward morphism complexes
are acyclic by \eqref{eq:trial-backward}. Replacing those complexes
by zero gives a quasi-equivalent upper triangular dg category,
whose diagonal dg algebras are the $E_v$.
Proposition~\ref{prop:trial-generation} makes that dg category
Morita equivalent to $A$. The smoothness theorem for dg gluings
\cite[Theorem 3.25]{Orl16} implies that every $E_v$ is smooth.
Now \cite[Proposition 7.4]{RS22} gives $E_v\simeq\Bbbk$.
Equation~\eqref{eq:trial-diagonal} proves (ii).

\noindent\underline{$(\mathrm{ii})\Rightarrow(\mathrm{iii})$.}
The brick condition and \eqref{eq:trial-diagonal} give the
degree-zero diagonal multiplicity. Thus (ii) strengthens
\eqref{eq:proposal-weak} to \eqref{eq:proposal-strong}.
Apply Proposition~\ref{prop:trial-strong}.

\noindent\underline{$(\mathrm{iii})\Rightarrow(\mathrm{i})$.}
Apply Corollary~\ref{cor:smooth}.
\end{proof}

\begin{center}\small
\begin{tabularx}{\linewidth}{@{}>{\raggedright\arraybackslash}p{.36\linewidth}>{\raggedright\arraybackslash}X@{}}
\toprule
Property under \eqref{eq:proposal-weak} & Result with
\eqref{eq:trial-brick} and condition~\ref{ax:filtration} \\\midrule
Classical specialisation & Yes; the negative degrees disappear.\\
Backward Hom in every degree & Yes; \eqref{eq:trial-backward}.\\
Bricks; positive self-Hom vanishing & Yes; \eqref{eq:trial-weak-self}.\\
Standards generate $\per A$ & Yes; \eqref{eq:trial-generation}.\\
Self-exceptionality; generation of $\D_{\fd}(A)$ & Not forced;
Example~\ref{ex:trial-dualnumbers}.\\
Smoothness as a consequence & No; the dg dual numbers are admitted.\\
Exact obstruction to smoothness & Negative diagonal cohomology;
Proposition~\ref{prop:trial-smooth-diagonal}.\\
\bottomrule
\end{tabularx}
\end{center}

\section{Adding smoothness as an assumption}
\label{sec:proposal-assumed-smooth}

For comparison, replace condition~\ref{ax:exceptional} by
\eqref{eq:trial-brick} and
\begin{equation*}\tag{C${}_{\mathrm{sm}}$}\label{eq:proposal-assumed-smooth}
 \text{\eqref{eq:proposal-weak} holds, and }
 A\in\per(A^{\op}\otimes_{\Bbbk}A).
\end{equation*}
Proposition~\ref{prop:trial-smooth-diagonal} proves that
\eqref{eq:proposal-assumed-smooth}, together with
condition~\ref{ax:filtration}, is equivalent to
Definition~\ref{def:shifted}. In particular, the attempted allowance
of negative diagonal cohomology disappears.

The algebra $\Bbbk[\varepsilon]/(\varepsilon^2)$ with
$|\varepsilon|=-1$ and $d=0$ is excluded precisely by the smoothness
assumption; it satisfies the other requirements by
Example~\ref{ex:trial-dualnumbers}. No non-smooth example satisfies
\eqref{eq:proposal-assumed-smooth}, by definition. The defect of this
proposal is axiomatic: smoothness is required rather than obtained
from the standard-module conditions.

\begin{center}\small
\begin{tabularx}{\linewidth}{@{}>{\raggedright\arraybackslash}p{.36\linewidth}>{\raggedright\arraybackslash}X@{}}
\toprule
Property under \eqref{eq:proposal-assumed-smooth} & Result with
\eqref{eq:trial-brick} and condition~\ref{ax:filtration} \\\midrule
Classical specialisation & Yes.\\
Backward orthogonality and self-exceptionality & Consequences of
Proposition~\ref{prop:trial-smooth-diagonal}.\\
Generation of $\per A=\D_{\fd}(A)$ & Yes.\\
Negative diagonal cohomology & Necessarily zero.\\
Non-smooth examples & Excluded by assumption.\\
Smoothness as a consequence of standard-module axioms & No;
smoothness is a separate hypothesis.\\
Distinct class from \eqref{eq:proposal-strong} & No;
the two formulations are equivalent under the filtration axiom.\\
\bottomrule
\end{tabularx}
\end{center}

\section{Projective tests with non-negative shifts}
\label{sec:proposal-positive}

Retain the brick requirement \eqref{eq:trial-brick} and the projective
presentations in condition~\ref{ax:filtration}. Consider replacing
condition~\ref{ax:exceptional} by
\begin{equation*}\tag{P+}\label{eq:proposal-positive}
 \Hom_{\D_{\fd}(A)}(P_w,\std_v[k])=0
 \qquad(v<_\sigma w,\ 0\leq k<m).
\end{equation*}
The ambient category in \eqref{eq:proposal-positive} is essential:
$\Hom_{\H_A^m}(P_w,\std_v[k])$ is defined only when
$\std_v[k]\in\H_A^m$. For every $v,w$ and $k\in\ZZ$, however,
\begin{equation}\label{eq:positive-projective-cohomology}
 \Hom_{\D_{\fd}(A)}(P_w,\std_v[k])
 \cong H^k(\std_v)e_w.
\end{equation}
Since $\std_v\in\H_A^m\subseteq\D_{\fd}(A)^{\leq0}$,
the terms in \eqref{eq:positive-projective-cohomology} vanish
automatically for $k>0$. Hence
\begin{equation}\label{eq:positive-reduction}
 \eqref{eq:proposal-positive}
 \quad\Longleftrightarrow\quad
 H^0(\std_v)e_w=0\qquad(v<_\sigma w).
\end{equation}
Restricting \eqref{eq:proposal-positive} to those $k$ for which
$\std_v[k]\in\H_A^m$ gives the same condition
\eqref{eq:positive-reduction}, since $k=0$ is always admissible.

For $m=1$, the standards are $H^0(A)$-modules, and
\eqref{eq:positive-reduction} says that the composition factors of
$\std_v$ have weights $u\leq_\sigma v$. Together with
\eqref{eq:trial-brick} and condition~\ref{ax:filtration}, these are
classical quasi-hereditary axioms. For arbitrary $m$,
Proposition~\ref{prop:trial-generation} gives
\begin{equation}\label{eq:positive-perfect-generation}
 \std_v\in\per A\quad(1\leq v\leq n),\qquad
 \thick_{\D_{\fd}(A)}(\std_1,\ldots,\std_n)=\per A.
\end{equation}
Indeed, descending induction in $\leq_\sigma$ proves perfection
from $\std_{\sigma(n)}\cong P_{\sigma(n)}$, while the projective
presentations express every $P_v$ in the thick subcategory generated
by the standards.

\begin{example}[Two failures of standard orthogonality]
\label{ex:positive-fork}
For $d\in\{1,2\}$, let
\begin{equation}\label{eq:positive-fork-algebra}
 A_d=\Bbbk\bigl(1\xrightarrow{a}3\xleftarrow{b}2\bigr),
 \qquad |a|=-d,\quad |b|=-1,\quad d_{A_d}=0,\quad m=d+1.
\end{equation}
Use right modules, so $e_1a=a=ae_3$ and $e_2b=b=be_3$.
Take $\sigma=\mathrm{id}$ and
\begin{equation}\label{eq:positive-fork-standards}
 (\std_1,\std_2,\std_3)=(P_1,S_2,P_3),\qquad
 (K_1,K_2,K_3)=(0,P_3[1],0).
\end{equation}
The presentations for vertices $1$ and $3$ are identities. For
vertex $2$, the short exact sequence of dg modules
\begin{equation}\label{eq:positive-fork-presentation}
 0\longrightarrow P_3[1]\xrightarrow{\iota_b}P_2
 \xrightarrow{p_2}S_2\longrightarrow0,
 \qquad \iota_b(e_3)=b,
\end{equation}
has all terms in $\H_{A_d}^{d+1}$; here $p_2$ is the quotient map.
Thus condition~\ref{ax:filtration} holds, with the single allowed
factor $K_2=\std_3[1]$. Each $\std_i$ is exceptional, and
$H^0(\std_i)=S_i$ for $i\in\{1,2,3\}$, so both
\eqref{eq:trial-brick} and \eqref{eq:proposal-positive} hold.

For every $t\in\ZZ$, the backward derived Hom spaces are
\begin{equation}\label{eq:positive-fork-hom}
 \begin{aligned}
 \Hom_{\D_{\fd}(A_d)}(\std_2,\std_1[t])
   &\cong\begin{cases}\Bbbk,&t=2-d,\\0,&t\neq2-d,\end{cases}\\
 \Hom_{\D_{\fd}(A_d)}(\std_3,\std_1[t])
   &\cong\begin{cases}\Bbbk,&t=-d,\\0,&t\neq-d,\end{cases}\\
 \Hom_{\D_{\fd}(A_d)}(\std_3,\std_2[t])&=0.
 \end{aligned}
\end{equation}
The second and third equalities follow from
\eqref{eq:positive-projective-cohomology}. For the first equality,
$\Hom_{\D_{\fd}(A_d)}(P_2,P_1[t])=0$ for every $t$, and the triangle
associated with \eqref{eq:positive-fork-presentation} gives
\begin{equation}\label{eq:positive-fork-calculation}
 \Hom_{\D_{\fd}(A_d)}(S_2,P_1[t])
 \cong\Hom_{\D_{\fd}(A_d)}(P_3[1],P_1[t-1])
 \cong H^{t-2}(P_1)e_3.
\end{equation}
Consequently:
\begin{enumerate}[label=\textup{(\roman*)}]
 \item $d=1$: backward degree-zero orthogonality holds, but
 $\Hom_{\D_{\fd}(A_1)}(\std_2,\std_1[1])\cong\Bbbk$;
 \item $d=2$: even backward degree-zero orthogonality fails,
 since $\Hom_{\D_{\fd}(A_2)}(\std_2,\std_1)\cong\Bbbk$.
\end{enumerate}
Both choices also violate backward negative-degree orthogonality
by \eqref{eq:positive-fork-hom}.
\end{example}

\begin{example}[Membership imposed only for ordered pairs]
\label{ex:positive-pair-membership}
Let
\begin{equation}\label{eq:positive-product-algebra}
 R=\Bbbk[\varepsilon]/(\varepsilon^2),\qquad
 |\varepsilon|=-1,\quad d_R=0,\qquad
 A=\Bbbk\times R,
\end{equation}
and take $m=2$, $\sigma=\mathrm{id}$, $\std_i=P_i$, and $K_i=0$
for $i\in\{1,2\}$. Each standard is a brick, and
\eqref{eq:proposal-positive} holds. Suppose additionally that
$\std_v[k]\in\H_A^2$ is required for the same pairs and shifts as
in \eqref{eq:proposal-positive}. The only pair is $(v,w)=(1,2)$,
and both $\std_1$ and $\std_1[1]$ belong to $\H_A^2$.
This pairwise membership requirement places no restriction on
$\std_2[1]$, which does not belong to $\H_A^2$.

The second direct factor of \eqref{eq:positive-product-algebra}
is the algebra in Example~\ref{ex:trial-dualnumbers}. Consequently,
$A$ is not homologically smooth,
$\Hom_{\D_{\fd}(A)}(\std_2,\std_2[-1])\cong\Bbbk$, and
\begin{equation}\label{eq:positive-product-not-full}
 \thick_{\D_{\fd}(A)}(\std_1,\std_2)
 =\per A\subsetneq\D_{\fd}(A).
\end{equation}
\end{example}

\begin{center}\small
\begin{tabularx}{\linewidth}{@{}>{\raggedright\arraybackslash}p{.36\linewidth}>{\raggedright\arraybackslash}X@{}}
\toprule
Property under \eqref{eq:proposal-positive} & Result with
\eqref{eq:trial-brick} and condition~\ref{ax:filtration} \\
\midrule
Classical specialisation at $m=1$
 & Yes: \eqref{eq:positive-reduction} is the classical support condition. \\
Generation of $\per A$
 & Yes, by \eqref{eq:positive-perfect-generation}. \\
Backward degree-zero orthogonality
 & Not forced: Example~\ref{ex:positive-fork}, $d=2$. \\
Backward positive-degree orthogonality
 & Not forced: Example~\ref{ex:positive-fork}, $d=1$. \\
Exceptionality of individual standards
 & Not forced: Example~\ref{ex:positive-pair-membership}. \\
Homological smoothness
 & Not forced: Example~\ref{ex:positive-pair-membership}. \\
Generation of $\D_{\fd}(A)$
 & Not forced: \eqref{eq:positive-product-not-full}. \\
\bottomrule
\end{tabularx}
\end{center}


\section{Requiring every shifted standard to remain in the extended category}
\label{sec:proposal-heart}

Retain \eqref{eq:trial-brick}, condition~\ref{ax:filtration}, and
\eqref{eq:proposal-positive}, and additionally require shift membership
for every standard, including the greatest weight:
\begin{equation}\label{eq:positive-all-membership}
 \std_v[k]\in\H_A^m
 \qquad(1\leq v\leq n,\ 0\leq k<m)
 \quad\Longleftrightarrow\quad
 \std_v\in\mod H^0(A)\qquad(1\leq v\leq n).
\end{equation}
Equation~\eqref{eq:positive-all-membership} is not part of
\eqref{eq:proposal-positive}. With
\eqref{eq:positive-all-membership},
\eqref{eq:positive-reduction} controls all cohomology of each
standard. Descending induction using
condition~\ref{ax:filtration} then gives
\eqref{eq:semiorthogonal}; the same projective presentations give
\begin{equation}\label{eq:positive-heart-self-hom}
 \Hom_{\D_{\fd}(A)}(\std_v,\std_v[t])
 \cong H^t(\std_v)e_v
 \qquad(1\leq v\leq n,\ t\in\ZZ).
\end{equation}
The heart condition in \eqref{eq:positive-all-membership} and the
brick requirement \eqref{eq:trial-brick} make each standard
exceptional. Thus imposing \eqref{eq:positive-all-membership}
restricts the standards to ordinary $H^0(A)$-modules and recovers
condition~\ref{ax:exceptional}.


\begin{example}[An excluded order on a graded arrow]
\label{ex:heart-excluded-order}
Let $A=\Bbbk(1\xrightarrow{a}2)$, with $e_1a=a=ae_2$,
$|a|=-1$, $d_A=0$, and $m=2$. For the order $2<_\sigma1$,
condition~\ref{ax:filtration} forces $K_1=0$ and $\std_1\cong P_1$.
Since $H^{-1}(P_1)=S_2$, we have $\std_1[1]\notin\H_A^2$.
No choice of standards for this order can satisfy
\eqref{eq:positive-all-membership}. The presentation
$\std_i=P_i$, $K_i=0$ satisfies \eqref{eq:proposal-strong}, so the
extra membership condition excludes an order admitted by the
exceptional-standard definition.

For the opposite order $1<_\sigma2$, the standards
$\std_1=S_1$, $\std_2=S_2$ and kernels $K_1=S_2[1]$, $K_2=0$
satisfy \eqref{eq:positive-all-membership},
\eqref{eq:proposal-positive}, and condition~\ref{ax:filtration}.
Thus this stronger condition restricts orders without requiring
$A$ itself to be concentrated in degree zero.
\end{example}

The dg dual numbers of Example~\ref{ex:trial-dualnumbers} are
excluded because the unique standard is $R$ and
$R[1]\notin\H_R^2$. No non-smooth algebra satisfies the requirements
of Section~\ref{sec:proposal-heart}: exceptionality follows from
\eqref{eq:positive-heart-self-hom}, and Corollary~\ref{cor:smooth}
applies.

\begin{center}\small
\begin{tabularx}{\linewidth}{@{}>{\raggedright\arraybackslash}p{.36\linewidth}>{\raggedright\arraybackslash}X@{}}
\toprule
Property & Result with \eqref{eq:positive-all-membership},
\eqref{eq:proposal-positive}, \eqref{eq:trial-brick}, and condition~\ref{ax:filtration}\\
\midrule
Classical specialisation & Yes; shift membership adds nothing for $m=1$.\\
Backward orthogonality; exceptional terms & Both are consequences.\\
Generation of $\per A=\D_{\fd}(A)$ & Yes.\\
Smoothness as a consequence & Yes; Corollary~\ref{cor:smooth}.\\
Negative cohomology of standards & Forbidden at every vertex.\\
Orders on graded acyclic path algebras & Some are excluded;
Example~\ref{ex:heart-excluded-order}.\\
Non-smooth examples & None satisfy these requirements.\\
\bottomrule
\end{tabularx}
\end{center}

\section{Projective tests detecting negative cohomology}
\label{sec:proposal-negative}

Retain the brick condition \eqref{eq:trial-brick}, the shifted-filtration
axiom~\ref{ax:filtration}, and $\std_v\in\H_A^m$ for every vertex $v$.
For $v,w\in\{1,\ldots,n\}$, consider the following
\defn{projective tests} $(\mathrm P^-)$:
\begin{equation*}\tag{P-}\label{eq:proposal-negative}
 \begin{aligned}
 \Hom_{\D_{\fd}(A)}(P_w,\std_v[-s])&=0
 &&(v<_\sigma w,\ 0\leq s<m),\\
 \Hom_{\D_{\fd}(A)}(P_v,\std_v[-s])&=0
 &&(1\leq s<m).
 \end{aligned}
\end{equation*}
Since
\begin{equation}\label{eq:proposal-projective-cohomology}
 \Hom_{\D_{\fd}(A)}(P_w,\std_v[-s])
 \cong H^{-s}(\std_v)e_w,
\end{equation}
the first line of \eqref{eq:proposal-negative} requires every
$H^t(\std_v)$ to have composition factors labelled $\leq_\sigma v$.
The second line excludes $S_v$ from $H^t(\std_v)$ for $t<0$.
Under the first line, Proposition~\ref{prop:trial-support} gives
\[
 H^0(\std_v)e_v
 \cong\Hom_{\D_{\fd}(A)}(\std_v,\std_v)
 \cong\Bbbk,
\]
where the last isomorphism uses \eqref{eq:trial-brick}.
Consequently \eqref{eq:proposal-negative} is equivalent to
\eqref{eq:proposal-strong} under the retained axioms.
Proposition~\ref{prop:trial-strong} therefore gives the exceptional
sequence condition, fullness in $\D_{\fd}(A)=\per A$, and smoothness
as consequences. For $m=1$, the second line of
\eqref{eq:proposal-negative} is empty, and the first line is an
ordinary Hom-vanishing condition for modules.

The tests \eqref{eq:proposal-negative} also admit a formulation using
only Hom spaces in $\H_A^m$. For every vertex $w$ and $0\leq s<m$, set
\begin{equation}\label{eq:proposal-truncated-projectives}
 U_{w,s}:=\tau^{\geq1-m}(P_w[s]),\qquad
 L_{w,s}:=\tau^{\leq-m}(P_w[s]).
\end{equation}
The amplitude of $P_w[s]$ implies $U_{w,s}\in\H_A^m$. The truncation
triangle is
\begin{equation}\label{eq:proposal-truncation-triangle}
 L_{w,s}\longrightarrow P_w[s]
 \xrightarrow{q_{w,s}}U_{w,s}\longrightarrow L_{w,s}[1].
\end{equation}
For $r\in\{0,1\}$, we have
$L_{w,s}[r]\in\D_{\fd}(A)^{\leq-m}$ and
$\std_v\in\D_{\fd}(A)^{\geq1-m}$. Thus
\[
 \Hom_{\D_{\fd}(A)}(L_{w,s}[r],\std_v)=0
 \qquad(r\in\{0,1\}).
\]
Applying $\Hom_{\D_{\fd}(A)}(-,\std_v)$ to
\eqref{eq:proposal-truncation-triangle} gives the natural isomorphisms
\begin{equation}\label{eq:proposal-intrinsic-projective-tests}
 \Hom_{\H_A^m}(U_{w,s},\std_v)
 \xrightarrow[\sim]{q_{w,s}^{*}}
 \Hom_{\D_{\fd}(A)}(P_w[s],\std_v)
 \cong H^{-s}(\std_v)e_w.
\end{equation}
Thus \eqref{eq:proposal-negative} can equivalently be imposed on
$\Hom_{\H_A^m}(U_{w,s},\std_v)$ with exactly the same ranges of
$v,w,s$. The truncations in \eqref{eq:proposal-truncated-projectives}
are necessary in this formulation: $P_w[s]$ need not belong to
$\H_A^m$ when $s>0$.

\begin{example}[An excluded one-vertex algebra]
\label{ex:proposal-negative-dual-numbers}
Let $m=2$, let
$R=\Bbbk[\varepsilon]/(\varepsilon^2)$ with
$|\varepsilon|=-1$ and $d_R=0$, and take $\std_1=P_1=R$ and $K_1=0$.
The brick condition and shifted-filtration axiom hold. The first line
of \eqref{eq:proposal-negative} is vacuous, whereas
\[
 U_{1,1}=\tau^{\geq-1}(R[1])\simeq\Bbbk[1],\qquad
 \Hom_{\H_R^2}(U_{1,1},R)
 \cong H^{-1}(R)\cong\Bbbk.
\]
Here $\Bbbk$ is the right $R$-module defined by the augmentation
$R\to\Bbbk$. Thus the second line of
\eqref{eq:proposal-negative} excludes $R$; imposing only the first
line would still admit $R$.
\end{example}

\begin{example}[Standards with negative cohomology]
\label{ex:proposal-negative-graded-arrow}
Let $A=\Bbbk Q$ for $Q=(1\xrightarrow{a}2)$, with $|a|=-1$ and
$d_A=0$. Take $m=2$, $2<_\sigma1$, and
$\std_1=P_1$, $\std_2=P_2$. Both presentation kernels are zero, and
\[
 H^0(\std_1)\cong S_1,\qquad
 H^{-1}(\std_1)\cong S_2,\qquad
 H^0(\std_2)\cong S_2.
\]
All other cohomology groups vanish. Therefore
\eqref{eq:proposal-negative} holds, and the standard sequence is
$(\std_2,\std_1)=(P_2,P_1)$. In particular, the projective tests permit
negative cohomology at strictly smaller vertices and do not require
the standards to lie in the heart $\mod H^0(A)$.
\end{example}

\begin{center}\small
\begin{tabularx}{\linewidth}{@{}>{\raggedright\arraybackslash}p{.36\linewidth}>{\raggedright\arraybackslash}X@{}}
\toprule
Property & Consequence of \eqref{eq:proposal-negative},
\eqref{eq:trial-brick}, and \ref{ax:filtration}\\
\midrule
Specialisation at $m=1$ & The classical quasi-hereditary definition.\\
Backward orthogonality & Vanishing in every degree.\\
Self-orthogonality & Every $\std_v$ is exceptional.\\
Fullness & $\boldsymbol\std^\sigma$ is full in
$\D_{\fd}(A)=\per A$.\\
Smoothness & A consequence of Proposition~\ref{prop:trial-strong}.\\
Negative standard cohomology & Allowed at vertices strictly smaller
than the label of the standard.\\
Non-smooth counterexamples & None, by
Proposition~\ref{prop:trial-strong}.\\
\bottomrule
\end{tabularx}
\end{center}


\section{Immediate problems}\label{sec:immediate}

\begin{enumerate}[label=\textbf{U\arabic*.},ref=U\arabic*]
\item\label{prob:comparison}
\textbf{Compare separate and combined amplitude bounds.}
Suppose a full exceptional sequence has presentations
\eqref{eq:presentation} with
$K_v\in\mathcal Q_v$, $\std_v\in\mathcal Q_v^\perp$,
$\std_v\in\H_A^m$, and $\costd_v\in\H_A^m[1-m]$.
The two amplitude conditions separately give
$\alpha_v\leq m-1$ and $\beta_v\leq m-1$.
Determine whether there are such examples with
$\alpha_v+\beta_v>m-1$ for some $v$.
An example would separate these conditions from
Definition~\ref{def:shifted} at the same width.
The degreewise obstruction is precisely the non-vanishing in
\eqref{eq:criterion} outside the allowed shifts.

\item\label{prob:recognition}
\textbf{Recognise the canonical standards from the dga.}
Given a basic proper connective dga $A$ and an order $\sigma$,
determine when every $P_v$ admits a triangle
$K_v\to P_v\to\std_v\to K_v[1]$ with
$K_v\in\mathcal Q_v$ and $\std_v\in\mathcal Q_v^\perp$,
where $\std_v$ is exceptional and lies in $\D_{\fd}(A)^{\leq0}$.
Seek criteria on the dg ideals $Ae_{>v}^\sigma A$ and their gluing
bimodules. Once the exceptional standard and left-dual objects
are constructed, Proposition~\ref{prop:criterion} tests the
filtration axiom, and Corollary~\ref{cor:width} determines the width
whenever one width exists.
\end{enumerate}

\section{Further problems}\label{sec:further}

\begin{enumerate}[label=\textbf{P\arabic*.}]
\item \textbf{Standardisation.}
Starting from an ordered exceptional family and allowed shifted
filtrations in an extriangulated category, determine the conditions
for a dg algebra to realise Definition~\ref{def:shifted}.
Specify the necessary enhancement and higher-composition data.
Compare with \cite{AT24}, checking its negative-extension hypotheses.

\item \textbf{Ringel duality.}
Determine when the categories generated by allowed shifted standards
and by the shifted costandards in \eqref{eq:injective} admit a common
tilting or silting object $T\in\D_{\fd}(A)$.
For the derived endomorphism dga $\REnd_A(T)$, determine connectivity,
the induced ordering, the width, and the effect of iterating the
construction.

\item \textbf{Homological bounds.}
Seek bounds depending on $m$ and $n$ for lengths of standard
resolutions and for
\[
 \sup\{t\in\ZZ:\Hom_{\D_{\fd}(A)}(S_u,S_v[t])\neq0
                \text{ for some }1\leq u,v\leq n\}.
\]

\item \textbf{Heredity ideals and degree-zero reduction.}
Relate \eqref{eq:localisation} to derived idempotent quotients and
chains of dg ideals. Determine the additional conditions on the
gluing data that encode \eqref{eq:kernel}, in comparison with the
classical recollement characterisation \cite{Kra17}.
Although \eqref{eq:support} identifies $H^0(\std_v)$ with the
ordinary standard module, Example~\ref{ex:differential} shows that
$H^0(A)$ need not be quasi-hereditary for the same ordering.
Characterise when quasi-heredity descends to $H^0(A)$.

\item \textbf{Changes of order and equivalence.}
Determine which mutations of $\boldsymbol\std^\sigma$ preserve
\eqref{eq:presentation} and \eqref{eq:kernel} at a fixed width.
Study how the standard objects and the least-width formula
\eqref{eq:minimal-width} change under silting mutation of the
projective generator and under changes of total order on the fixed
vertex labels.
\end{enumerate}

\begin{thebibliography}{Kra17}
\bibitem[AT24]{AT24} T. Adachi and M. Tsukamoto,
\emph{Mixed standardization and Ringel duality},
J. Algebra \textbf{641} (2024), 546--586.
\href{https://arxiv.org/abs/2202.06057}{arXiv:2202.06057}.

\bibitem[BB26]{BB26} A. Bodzenta and A. Bondal,
\emph{Highest weight categories via pairs of dual exceptional sequences},
2026. \href{https://arxiv.org/abs/2601.22004}{arXiv:2601.22004}.

\bibitem[Kra17]{Kra17} H. Krause,
\emph{Highest weight categories and recollements},
Ann. Inst. Fourier \textbf{67} (2017), 2679--2701.
\href{https://arxiv.org/abs/1506.01485}{arXiv:1506.01485}.

\bibitem[MP25]{MP25} N. Mochizuki and M. Plogmann,
\emph{On the Auslander--Reiten theory for extended hearts of proper
connective dg algebras}, 2025.
\href{https://arxiv.org/abs/2505.16560}{arXiv:2505.16560}.

\bibitem[Orl16]{Orl16} D. Orlov,
\emph{Smooth and proper noncommutative schemes and gluing of DG categories},
Adv. Math. \textbf{302} (2016), 59--105.
\href{https://arxiv.org/abs/1402.7364}{arXiv:1402.7364}.

\bibitem[RS22]{RS22} T. Raedschelders and G. Stevenson,
\emph{Proper connective differential graded algebras and their geometric
realizations}, European J. Math. \textbf{8} (2022), S574--S598.
\href{https://arxiv.org/abs/1903.02849}{arXiv:1903.02849}.

\end{thebibliography}
\end{document}
